\section{Methodology}
\label{sec:method}

\subsection{AR Plan Generation}

The \ar{} plan for each state is generated by the \texttt{redist} binary
at the minimum-edge-cut seed found by \cs{} at threshold $T = 600$
\citep{dellaLibera2026b16}.  Full details of the algorithm are in
\citet{dellaLibera2026b11}.  Briefly: the seat count $k$ is factored into
primes; each prime factor $p$ generates a $p$-way minimum-cut partition of
the current subgraph; the hierarchy continues until every leaf contains
exactly one district.  Where $k$ is prime (Pennsylvania, $k = 17$), a
single $k$-way partition is used.

\subsection{Ensemble Comparison}

For states with published ensemble results, we extract:
\begin{itemize}
  \item The empirical distribution of mean Polsby-Popper scores
    $\bar\pp \in (0, 1)^N$;
  \item The empirical distribution of Democratic seat counts
    $S_D \in \{0, \ldots, k\}^N$;
  \item The empirical distribution of minority-majority district counts
    $\mathrm{MM} \in \{0, \ldots, k\}^N$.
\end{itemize}

The \ar{} plan's percentile on each statistic is computed via
equation~(1) from G.0.

\subsection{Partisan Baseline}

All partisan results use the 2020 presidential vote as the electoral baseline.
Democratic percentage by census tract is computed from the \texttt{redist}
analysis output (the \texttt{redist analyze} command with
\texttt{--types partisan}).  For the ensemble comparison, we use the same
baseline as the published ensemble where possible.

\subsection{Uncertainty Quantification}

For literature-bound estimates (TX, CA), we report a 90\% confidence interval
for the percentile position based on normal approximation to the ensemble
distribution:

\begin{equation}
  \hat\pi \pm 1.645 \cdot \sqrt{\hat\pi(1-\hat\pi)/N}.
\end{equation}

For direct comparisons (NC, WI, GA, PA), the percentile estimate has
negligible sampling error given $N \geq 10{,}000$.
